\documentclass[12pt]{article}
\usepackage{amsmath,amssymb}
\begin{document}
\section*{Research notes}

Problem: uniqueness of invariant probability for Dirichlet stochastic heat
equation on $[0,1]$ with additive space-time white noise and distributional
drift $\delta_0(u)$, under the ABLM regularized/mild interpretation and
regularisation by $G_{1/n}$.

Key interpretation:
\[
\delta_0(r)=-\frac12(-2{\bf1}_{r>0})',
\]
so the equation is a singular gradient SPDE with bounded potential
$f(r)=-2H_+(r)$.  For smooth approximation $B_n'=G_{1/n}$ with
$B_n(r)=\int_{-\infty}^rG_{1/n}(s)\,ds$, the invariant measure of the
regularized gradient equation is
\[
\pi_n(dv)=Z_n^{-1}\exp\left(2\int_0^1 B_n(v(x))\,dx\right)\gamma(dv),
\]
where $\gamma$ is Brownian bridge law (the invariant measure of the
Dirichlet OU/string process).  Since $0\le B_n\le1$ and $B_n\to H_+$
away from zero, while a Brownian bridge spends zero Lebesgue time at
zero,
\[
\pi_n\to
\pi(dv)=Z^{-1}\exp\left(2\int_0^1{\bf1}_{v(x)>0}\,dx\right)\gamma(dv)
\]
in total variation.  Densities $d\pi_n/d\gamma$ and $d\pi/d\gamma$ are
bounded above and below by deterministic positive constants.

Round-1 robust result obtained: assuming pointwise weak convergence of
regularized semigroups on bounded continuous observables,
\[
P_t^n\Phi(x)\to P_t\Phi(x),\qquad x\in C_0,\ t>0,
\]
one can pass reversibility/invariance and a uniform spectral gap from
$P^n$ to the actual ABLM semigroup $P$ without invoking
Bounebache--Zambotti:
\[
\int gP_tf\,d\pi=\int fP_tg\,d\pi,\qquad
\operatorname{Var}_\pi(P_tf)\le e^{-2ct}\operatorname{Var}_\pi(f).
\]
Key estimates:
\[
\operatorname{Var}_{\pi_n}(F)\le C\mathcal E_n(F,F),\qquad
\mathcal E_n(F,F)=\frac12\int\|DF\|_{L^2}^2\,d\pi_n,
\]
with $C$ uniform in $n$, by Gaussian Poincare for $\gamma$ plus bounded
density perturbation.  Then pass the $L^2$ contraction/spectral gap to
the limit using total variation convergence $\pi_n\to\pi$ and bounded
pointwise convergence $P_t^nf\to P_tf$.

Abstract uniqueness mechanism:
If $\eta$ is invariant and the actual resolvent satisfies
\[
R_\lambda^P(x,\cdot)\ll\pi\quad\hbox{for every }x,
\]
then $\eta=\lambda\eta R_\lambda^P\ll\pi$.  With $h=d\eta/d\pi$,
symmetry gives $P_th=h$ in $L^1(\pi)$.  For $h_a=h\wedge a$, concavity
gives $h_a\ge P_th_a$ and equal integrals force equality.  Since
$h_a\in L^2$, the spectral gap forces $h_a$ constant.  Varying $a$
implies $h$ constant, hence $\eta=\pi$.

Central unresolved issue:
Need pointwise resolvent absolute continuity for the actual ABLM
semigroup, not only for the $L^2(\pi)$ Dirichlet-form version and not
only quasi-everywhere.  BZ Proposition 2.5 gives an $L^2$ resolvent kernel
$\rho_\lambda(x,dy)\ll\nu(dy)$ for all $x$ (setting it to zero on an
exceptional set) and transition absolute continuity for all
non-exceptional initial conditions.  BZ Proposition 3.1 constructs a
weak local-time solution for quasi-every initial condition.  ABLM
uniqueness plus the 2025 weak/mild equivalence paper plausibly identify
the BZ process with the ABLM mild process for quasi-every initial point,
but this still does not prove that arbitrary invariant measures cannot
charge the exceptional set, nor that $R_\lambda^P(x,\cdot)\ll\pi$ for
$x$ in that set.

Potential routes to close the gap:
1. Uniform heat-kernel/capacity estimates for $P_t^n$: prove
   $P_t^n(x,A)\le C_{t,x}\,\mathrm{Cap}_n(A)^\alpha$ or
   $\|dP_t^n(x,\cdot)/d\pi_n\|_{L^2(\pi_n)}\le C_{t,x}$ with constants
   stable as $n\to\infty$ (perhaps only for $t$ large).  This would pass
   AC to the limit and imply uniqueness via convergence to equilibrium.
2. Strict version of BZ process: prove that the ABLM process is the
   strict all-starting-points Hunt process associated with the BZ
   Dirichlet form, and that BZ's resolvent kernel represents its
   resolvent at every deterministic $x$.
3. Entrance from exceptional sets: let $N$ be the BZ exceptional set.
   Show $P_s^P(x,N)=0$ for every $x$ and $s>0$.  Then q.e. resolvent AC
   plus the Markov property gives all-$x$ resolvent AC:
   $R_\lambda(x,A)\le\int_0^s e^{-\lambda t}P_t(x,A)dt$ for
   $\pi(A)=0$, then let $s\downarrow0$.
4. Prove uniqueness of invariant measures without pointwise AC by showing
   every invariant stationary solution has enough finite-energy/Sobolev
   regularity not to charge zero-capacity sets.  Stationarity and the
   nondegenerate additive noise imply one-dimensional projections have no
   atoms by the occupation formula, but this is far weaker than avoiding
   all Gaussian-capacity-zero sets.

Useful literature facts:
- BZ use $\nu(dx)\propto \exp(-\int f(x_r)dr)\mu(dx)$.  For
  $f=-2H_+$ this is exactly $\pi$.
- BZ equation has singular term
  $-\frac12\int f(da)\partial_\theta\ell^a_{t,\theta}$; for
  $f=-2H_+$ this is $\partial_\theta\ell^0$, formally
  $\delta_0(u(t,\theta))$.
- ABLM 2024 proves strong existence/uniqueness for $b=\kappa\delta_0$
  and convergence of smooth approximations; it explicitly did not treat
  Dirichlet boundary conditions, but the problem assumes the extension.
- ABLM 2025 ``Analytically weak and mild solutions...'' proves weak
  regularized and mild regularized solution notions coincide in the
  known existence range; may be useful for local-time-to-mild
  identification, though it still does not by itself remove the
  exceptional-set issue.

\end{document}
